\documentclass[12pt]{article}
\usepackage[T1]{fontenc}
\usepackage[utf8]{inputenc}
\usepackage{lmodern}
\usepackage{textcomp}
\usepackage{enumitem}
\usepackage{geometry}
\geometry{margin=1in}
\begin{document}
\begin{center}
Answer Key - Health Sciences - K-12 Level Assessment 20 \
\small (Answers are ordered exactly as in the question paper.)
\end{center}

\section*{Multiple Choice}
\begin{enumerate}[label=\arabic*.]
\item \textbf{Answer:} D
\\ \textit{Options:} A) family history; B) being overweight; C) high intake of dietary fat; D) All of the options listed are correct
\\ \textit{Note:} The correct answer is "All of the options listed are correct" because type 2 diabetes mellitus is influenced by a combination of various risk factors such as genetics, lifestyle choices, obesity, and age, making it essential to consider all contributing elements collectively.
\item \textbf{Answer:} D
\\ \textit{Options:} A) Lipids; B) Carbohydrates; C) Nucleotides; D) Proteins
\\ \textit{Note:} Histones are proteins that play a crucial role in packaging and organizing DNA into structural units called nucleosomes, which are essential for DNA compaction in cells.
\item \textbf{Answer:} B
\\ \textit{Options:} A) The Summit's goal of cutting the number of hungry people in half by 2015 was achieved.; B) The number decreased, but not by nearly enough to meet the Summit's goal.; C) The number increased slightly.; D) Because of rising food prices, the number increased dramatically.
\\ \textit{Note:} The number of food insecure people has decreased since the 1996 World Food Summit, but the reduction has been insufficient to achieve the Summit's goal of reducing hunger significantly, highlighting ongoing challenges in global food security.
\item \textbf{Answer:} C
\\ \textit{Options:} A) Fish and fish products; B) Pulses; C) Meat and meat products; D) Milk and milk products
\\ \textit{Note:} Meat and meat products are the main source of protein in the British diet because they provide a high-quality, complete protein that contains all essential amino acids needed for growth and maintenance of body tissues.
\item \textbf{Answer:} C
\\ \textit{Options:} A) A baked apple; B) A raw apple; C) A raw potato; D) Apple juice
\\ \textit{Note:} A raw potato has a lower glycaemic index compared to other foods because it contains resistant starch, which slows down digestion and the release of glucose into the bloodstream.
\end{enumerate}

\section*{True / False}
\begin{enumerate}[label=\arabic*.]
\setcounter{enumi}{5}
\item \textbf{Answer:} False
\\ \textit{Options:} A) True; B) False
\\ \textit{Note:} The doubly labelled water technique primarily measures total energy expenditure rather than directly assessing food intake, particularly undervaluing consumption of low energy content foods which may not significantly influence metabolic water production.
\item \textbf{Answer:} False
\\ \textit{Options:} A) True; B) False
\\ \textit{Note:} The statement is false because patients with Bulimia Nervosa often have a body mass index (BMI) that is within the normal range or may even be overweight, as they engage in cycles of binge eating followed by compensatory behaviors rather than consistently underweight characteristics.
\item \textbf{Answer:} False
\\ \textit{Options:} A) True; B) False
\\ \textit{Note:} Lactobacilli contribute to dental caries development, but Streptococcus mutans is primarily responsible for the initial stages of tooth decay.
\item \textbf{Answer:} False
\\ \textit{Options:} A) True; B) False
\\ \textit{Note:} Demineralisation of tooth enamel occurs when the pH level falls below 5.5, not at 6.0, because a pH of 6.0 is still above the critical threshold needed for enamel to remain intact.
\item \textbf{Answer:} True
\\ \textit{Options:} A) True; B) False
\\ \textit{Note:} Total urinary nitrogen excretion can accurately determine the daily amount of protein oxidized in the body because nitrogen is a key component of amino acids, which are the building blocks of proteins, and measuring its excretion provides a reliable estimate of protein catabolism.
\end{enumerate}

\section*{Long Form Response}
\begin{enumerate}[label=\arabic*.]
\setcounter{enumi}{10}
\item \textbf{Answer:} Consuming vitamin C and other food supplements in excess can lead to various health issues. For example, high doses of vitamin C can cause gastrointestinal problems such as diarrhea, nausea, and stomach cramps. Additionally, excessive intake of certain supplements may interfere with the absorption of other essential nutrients, potentially leading to imbalances in the body. To ensure safe consumption, it is recommended to follow the recommended dietary allowances (RDAs) and consult with a healthcare professional before taking supplements. A balanced diet rich in fruits and vegetables typically provides sufficient vitamins without the risks associated with high doses.
\\ \textit{Note:} correct because it highlights the negative health effects of excessive vitamin C consumption, such as gastrointestinal issues and nutrient imbalances, while also emphasizing the importance of following recommended dietary allowances for safe supplement intake.
\item \textbf{Answer:} During digestion, water transport between the intestinal lumen and the bloodstream is crucial for maintaining proper osmolality and supporting nutrient absorption. As food, particularly starch and protein, is broken down in the intestines, water flows into the gut from the mucosa to help dilute the contents. This process reduces the luminal osmolality, which is the concentration of solutes in the intestinal fluid, making it easier for nutrients to be absorbed through the intestinal walls into the bloodstream. By ensuring that the osmotic balance is maintained, the body can efficiently absorb not only water but also essential nutrients, ultimately supporting overall health and digestion.
\\ \textit{Note:} correct because it accurately describes how water transport into the intestinal lumen aids in diluting the contents, thereby reducing osmolality and facilitating the absorption of nutrients into the bloodstream during digestion.
\end{enumerate}

\vfill
\noindent\textit{End of Answer Key}
\end{document}